\documentclass[11pt]{article}

\usepackage[margin=1in]{geometry}
\usepackage{amsmath,amssymb,amsthm}
\usepackage{booktabs}
\usepackage{array}
\usepackage{microtype}
\usepackage{xurl}
\usepackage{pgfplots}
\pgfplotsset{compat=1.18}
\setlength{\emergencystretch}{1em}
\usepackage{enumitem}
\usepackage{natbib}
\usepackage[hidelinks]{hyperref}

\newtheorem{proposition}{Proposition}
\newtheorem{assumption}{Assumption}
\newtheorem{lemma}{Lemma}
\newtheorem{corollary}{Corollary}
\newcommand{\R}{\mathbb{R}}
\newcommand{\dd}{\mathrm{d}}
\newcommand{\KDM}{\ensuremath{\mathrm{KDM}}}
\newcommand{\LEDH}{\ensuremath{\mathrm{LEDH}}}
\newcommand{\GenUT}{\ensuremath{\mathrm{GenUT}}}
\newcommand{\OT}{\ensuremath{\mathrm{OT}}}


\title{Likelihood Scores for Degenerate State-Space Models\\
Disturbance Proposals, Analytical Recursions, and Exact Control Variates}
\author{BayesFilter research note}
\date{12 September 2026}
\begin{document}
\maketitle
\begin{abstract}
When transition noise has lower dimension than the state, an ordinary
transition-density smoothing recursion can be undefined. We derive a score
construction in the disturbances that generate the state. Explicit
assumptions, propositions, and proofs establish the complete-path score,
unbiased particle likelihood and derivative estimates, an analytically
recursive implementation, and exactly centered reference control variates.
Conditional integration and a moving-support Gaussian example clarify the
mechanism. These results preserve the original model, but do not establish
an unbiased finite-particle normalized score or low variance at long horizons.
\end{abstract}
\section{The disturbance-coordinate proposal}

\label{sec:disturbance-proposal}

The support obstruction changes the coordinates in which the proposal must be
defined. Instead of proposing an ambient state and then asking for the density
of a singular transition, propose the random variables that generate the
state. This is the disturbance-state-space formulation of
\citet[Sections 1--2]{murray2018disturbance}. It is also the only formulation
used in the propositions below.

Let $v_0\in\mathsf V_0$ be an initial disturbance and let
$u_t\in\mathsf U_t$ be the innovation at time $t$. The sets
$\mathsf V_0$ and $\mathsf U_t$ carry fixed reference measures
$\lambda_0$ and $\lambda_t$; they need not have the same dimension as the
state. For a parameter $\theta$ in an open set $\Theta\subset\R^p$, define
\begin{align}
 x_0&=F_{0,\theta}(v_0), &
 x_t&=F_{t,\theta}(x_{t-1},u_t), \quad t=1,\ldots,T, 
 \label{eq:disturbance-state-map}\\
 y_t\mid x_t&\sim g_{t,\theta}(\,\cdot\mid x_t), &
 v_0&\sim\rho_{0,\theta}, \quad
 u_t\mid x_{t-1}\sim r_{t,\theta}(\,\cdot\mid x_{t-1}).
 \label{eq:disturbance-model}
\end{align}
The transition may be deterministic and rank deficient in $x_t$; no density
with respect to ambient Lebesgue measure on $\R^{d_x}$ is assumed. The
densities $\rho_{0,\theta}$, $r_{t,\theta}$, and $g_{t,\theta}$ are with
respect to their displayed fixed reference measures; the observation
reference measure is also independent of $\theta$. We assume that the
innovations have the displayed conditional Markov law given the past and
that observations are conditionally independent given the state path and
$\theta$. These are model assumptions. If shocks are independent of the state
and parameter, $r_{t,\theta}$ reduces to a fixed $r_t$.

\begin{assumption}[Disturbance regularity]
\label{ass:disturbance-regularity}
For $\lambda=\lambda_0\otimes\lambda_1\otimes\cdots\otimes\lambda_T$,
define $z=(v_0,u_{1:T})$ and let $x_{0:T}(z,\theta)$ be generated by
\eqref{eq:disturbance-state-map}. There is a neighbourhood $B$ of the
parameter value under study such that:
\begin{enumerate}[label=(\roman*)]
  \item on a common set of full $\lambda$ measure the state maps are
        continuously differentiable in their parameter and state arguments
        along the generated paths, and the displayed density factors are
        positive and continuously differentiable there. Their common
        disturbance support is independent of $\theta\in B$;
  \item the disturbance-coordinate integrand
        \begin{equation}
        \pi_\theta(z;y)=\rho_{0,\theta}(v_0)
        \prod_{t=1}^T r_{t,\theta}(u_t\mid x_{t-1})
        g_{t,\theta}(y_t\mid x_t)
        \label{eq:disturbance-joint-integrand}
        \end{equation}
        has a positive finite integral $Z(\theta;y)$ for $\theta\in B$; and
  \item each component of $\nabla_\theta\pi_\theta(z;y)$ is dominated on $B$ by an
        integrable envelope. The same conditions hold for each time prefix
        used below.
\end{enumerate}
\end{assumption}

The fixed-coordinate condition is the important one. A state support that
moves with $\theta$ is harmless when it is generated by the same fixed
disturbance coordinate $z$; differentiating an ambient state density is not
needed. For a parameter-dependent shock support, first introduce a fixed-base
variable and put its parameter dependence into $F_{t,\theta}$ and its density.

Use column gradients for scalar functions and let $\dot x_t$ be the
$d_x\times p$ Jacobian of $x_t$ with respect to $\theta$, holding the
disturbances fixed. Define this total tangent recursively by
\begin{align}
 \dot x_0&=\partial_\theta F_{0,\theta}(v_0),\\
 \dot x_t&=\partial_\theta F_{t,\theta}(x_{t-1},u_t)
       +J_xF_{t,\theta}(x_{t-1},u_t)\dot x_{t-1}.
 \label{eq:disturbance-tangent}
\end{align}
For any state-dependent shock density, use the total derivative
\begin{equation}
 D_\theta\log r_{t,\theta}(u_t\mid x_{t-1})
 =\partial_\theta\log r_{t,\theta}(u_t\mid x_{t-1})
  +\dot x_{t-1}^{\mathsf T}\nabla_x\log r_{t,\theta}(u_t\mid x_{t-1}).
 \label{eq:disturbance-shock-total-derivative}
\end{equation}
The complete disturbance-path score is
\begin{align}
 h_\theta(z;y)
 &=\nabla_\theta\log\rho_{0,\theta}(v_0)
   +\sum_{t=1}^T D_\theta\log r_{t,\theta}(u_t\mid x_{t-1})
 \nonumber\\[-2pt]
 &\quad+\sum_{t=1}^T\left[
       \partial_\theta\log g_{t,\theta}(y_t\mid x_t)
       +\dot x_t^{\mathsf T}\nabla_x\log g_{t,\theta}(y_t\mid x_t)\right].
 \label{eq:disturbance-complete-score}
\end{align}

\begin{proposition}[Disturbance-coordinate Fisher identity]
\label{prop:disturbance-fisher}
Under Assumption~\ref{ass:disturbance-regularity}, the observed-data
likelihood and score are
\begin{align}
 Z(\theta;y)&=\int \pi_\theta(z;y)\,\dd\lambda(z),\\
 \nabla_\theta\log Z(\theta;y)
 &=\mathbb E_\theta\!\left[h_\theta(\mathsf Z;y)\mid y_{1:T}\right],
 \label{eq:disturbance-score-identity}
\end{align}
where the conditional law in the second line has density
$\pi_\theta(z;y)/Z(\theta;y)$ with respect to $\lambda$.
\end{proposition}

\begin{proof}
The recursion in \eqref{eq:disturbance-state-map} makes
$\pi_\theta(z;y)$ a function on the fixed measure space
$(\mathsf V_0\times\mathsf U_{1:T},\lambda)$. By the envelope condition,
differentiation may pass through the integral, so
\begin{equation}
 \nabla_\theta Z(\theta;y)=\int \nabla_\theta\pi_\theta(z;y)\,\dd\lambda(z).
 \label{eq:disturbance-differentiate-integral}
\end{equation}
At points where the factors are positive, the product rule gives
$\nabla_\theta\pi_\theta=\pi_\theta h_\theta$. The chain rule through the
state recursion gives \eqref{eq:disturbance-tangent}; applying it to each
shock and observation factor gives \eqref{eq:disturbance-complete-score}.
Substitution into \eqref{eq:disturbance-differentiate-integral}, followed by
division by the positive $Z(\theta;y)$, gives
\begin{equation}
 \nabla_\theta\log Z(\theta;y)
 =\int h_\theta(z;y)\frac{\pi_\theta(z;y)}{Z(\theta;y)}\,\dd\lambda(z),
\end{equation}
which is the asserted conditional expectation. No transition density in the
ambient state coordinates appears in the argument.
\end{proof}

The initial term in \eqref{eq:disturbance-complete-score} is not a cosmetic
addition. If the model starts in its stationary distribution and that law
depends on $\theta$, the stationary derivative belongs in
$\nabla_\theta\log\rho_{0,\theta}$. If instead a parameter-independent
disturbance generates the initial state through $F_{0,\theta}$, its effect
appears in $\dot x_0$ and in every later observation term.

\begin{proposition}[Support-valid sequential proposal]
\label{prop:disturbance-proposal}
Under Assumption~\ref{ass:disturbance-regularity}, suppose an initial
proposal $q_{0,\theta}(v_0)$ and, for each $t$, a proposal
$q_{t,\theta}(u_t\mid x_{t-1},y_t)$ are positive wherever the corresponding
target factors in \eqref{eq:disturbance-joint-integrand} are positive. For
sequential importance sampling without resampling, draw from
\begin{equation}
 Q_\theta(z)=q_{0,\theta}(v_0)\prod_{t=1}^T
 q_{t,\theta}(u_t\mid x_{t-1},y_t),
 \label{eq:disturbance-proposal-law}
\end{equation}
and set, for a draw $\mathsf Z\sim Q_\theta$,
\begin{align}
 R_\theta(z)&=\frac{\pi_\theta(z;y)}{Q_\theta(z)},\\
 \widehat Z&=R_\theta(\mathsf Z), &
 \widehat D&=R_\theta(\mathsf Z)h_\theta(\mathsf Z;y).
 \label{eq:disturbance-importance-pair}
\end{align}
The differentiability and integrable derivative bound in
Assumption~\ref{ass:disturbance-regularity} imply that $Z$ is differentiable
and that $\pi_\theta h_\theta$ is integrable. Consequently,
\begin{equation}
 \mathbb E_{Q_\theta}[\widehat Z]=Z(\theta;y),\qquad
 \mathbb E_{Q_\theta}[\widehat D]=\nabla_\theta Z(\theta;y).
 \label{eq:disturbance-importance-identities}
\end{equation}
Where $\widehat Z>0$, the one-draw ratio is
$h_\theta(\mathsf Z;y)$. Its sampling law here is $Q_\theta$, so it is
generally not an unbiased model-score estimator. If $Q_\theta$ is the exact
disturbance posterior, even one such draw is unbiased by
Proposition~\ref{prop:disturbance-fisher}; an arbitrary proposal has no such
guarantee.
\end{proposition}

\begin{proof}
The support condition makes $R_\theta$ well defined $Q_\theta$-almost
everywhere. Cancellation of $Q_\theta$ gives
\begin{align*}
 \mathbb E_{Q_\theta}[R_\theta]
 &=\int \pi_\theta(z;y)\,\dd\lambda(z)=Z(\theta;y),\\
 \mathbb E_{Q_\theta}[R_\theta h_\theta]
 &=\int \pi_\theta(z;y)h_\theta(z;y)\,\dd\lambda(z)
 =\nabla_\theta Z(\theta;y),
\end{align*}
where the last equality is Proposition~\ref{prop:disturbance-fisher} before
division by $Z$. The proposal density may depend on $\theta$: it cancels in
the pointwise importance ratio. A derivative of $Q_\theta$ is needed only if
one differentiates the proposal expectation itself, rather than using the
score identity to estimate $\nabla_\theta Z$.
\end{proof}

For a particle filter, the same ratio is applied one step at a time. Let
$\alpha_t(a\mid\mathcal F_{t-1})$ be an ancestor proposal, positive whenever
$W_{t-1}^a>0$, and draw $u_t^i$ from $q_{t,\theta}$ after drawing
$a_t^i$. The support-valid auxiliary particle weight is
\begin{equation}
 \omega_t^i=
 \frac{W_{t-1}^{a_t^i}}{\alpha_t(a_t^i\mid\mathcal F_{t-1})}
 \frac{r_{t,\theta}(u_t^i\mid x_{t-1}^{a_t^i})
       g_{t,\theta}(y_t\mid x_t^i)}
      {q_{t,\theta}(u_t^i\mid x_{t-1}^{a_t^i},y_t)},
 \qquad x_t^i=F_{t,\theta}(x_{t-1}^{a_t^i},u_t^i).
 \label{eq:disturbance-apf-weight}
\end{equation}
If $\alpha_t=W_{t-1}$, the first ratio is one. A lookahead or guided
ancestor law changes $\alpha_t$ but not the numerator. This is the exact
place where a Younis or UKF construction may enter: it supplies
$\alpha_t$ or $q_t$ while the disturbance density and observation likelihood
remain those of the declared model.

\begin{proposition}[Conditional APF normalizer identity]
\label{prop:disturbance-apf-normalizer}
Conditionally on $\mathcal F_{t-1}$, assume that $\alpha_t$ and $q_{t,\theta}$
are fixed before the ancestor and innovation draws and that the ratios in
\eqref{eq:disturbance-apf-weight} are integrable. Then, for every particle,
\begin{equation}
 \mathbb E[\omega_t^i\mid\mathcal F_{t-1}]
 =\sum_{a=1}^N W_{t-1}^a
   \int r_{t,\theta}(u\mid x_{t-1}^a)
        g_{t,\theta}(y_t\mid F_{t,\theta}(x_{t-1}^a,u))\,\dd\lambda_t(u).
 \label{eq:disturbance-apf-conditional-normalizer}
\end{equation}
If the incoming empirical measure is replaced by the exact filtering measure,
the right side is the exact predictive likelihood.
\end{proposition}

\begin{proof}
Condition on an ancestor value $a$. The proposal density cancels in the
innovation integral:
\begin{align*}
 \mathbb E[\omega_t^i\mid\mathcal F_{t-1},a_t^i=a]
 &=\frac{W_{t-1}^a}{\alpha_t(a\mid\mathcal F_{t-1})}
   \int r_{t,\theta}(u\mid x_{t-1}^a)
        g_{t,\theta}(y_t\mid F_{t,\theta}(x_{t-1}^a,u))\,\dd\lambda_t(u).
\end{align*}
Taking expectation over $a_t^i\sim\alpha_t$ yields
\eqref{eq:disturbance-apf-conditional-normalizer}. Integrating the same
innovation likelihood over the exact incoming filtering measure gives the
predictive likelihood by the law of total probability.
\end{proof}

The one-step cancellation alone is not a proof about a complete filtering
run. To supply that proof, retain the disturbance history
$z_{0:t}=(v_0,u_{1:t})$ mathematically and write
$\gamma_t(\varphi)=\int\pi_{\theta,t}\varphi\,\dd\lambda_{0:t}$, where
$\pi_{\theta,t}$ is \eqref{eq:disturbance-joint-integrand} stopped at $t$.
There is no observation at time zero, so $\pi_{\theta,0}=\rho_{0,\theta}$
and $\gamma_0(1)=1$. The extension kernel for a history test function is
\begin{equation}
 (K_{t,\theta}\varphi)(z_{0:t-1})
 =\int r_{t,\theta}(u\mid x_{t-1})
       g_{t,\theta}(y_t\mid F_{t,\theta}(x_{t-1},u))
       \varphi(z_{0:t-1},u)\,\dd\lambda_t(u).
 \label{eq:disturbance-history-kernel}
\end{equation}
Direct integration of the last disturbance gives
$\gamma_t(\varphi)=\gamma_{t-1}(K_{t,\theta}\varphi)$.

\begin{proposition}[Unbiased particle likelihood and derivative]
\label{prop:disturbance-particle-induction}
Fix $N\geq1$ and $\theta$. Initially draw $v_0^i$ independently from
$q_{0,\theta}$ and set
$\omega_0^i=\rho_{0,\theta}(v_0^i)/q_{0,\theta}(v_0^i)$.
At each subsequent time, draw ancestors conditionally independently from
$\alpha_t$ and innovations from the proposals in
\eqref{eq:disturbance-apf-weight}. Suppose all weights are positive and
finite almost surely, all proposal support conditions hold, and the test
functions below are absolutely integrable under $\gamma_t$. Define
\begin{equation}
 \overline\omega_t=\frac1N\sum_i\omega_t^i,\qquad
 W_t^i=\frac{\omega_t^i}{N\overline\omega_t},\qquad
 \widehat Z_t=\prod_{s=0}^t\overline\omega_s,\qquad
 \gamma_t^N(\varphi)=\widehat Z_t\sum_iW_t^i\varphi(z_{0:t}^i).
 \label{eq:disturbance-particle-measure}
\end{equation}
For every deterministic test function $\varphi$ with the stated
integrability, including one that depends on the fixed $\theta$,
\begin{equation}
 \mathbb E[\gamma_t^N(\varphi)]=\gamma_t(\varphi).
 \label{eq:disturbance-particle-unbiasedness}
\end{equation}
In particular, with $h_{\theta,t}=\nabla_\theta\log\pi_{\theta,t}$,
\begin{equation}
 \widehat D_t=\widehat Z_t\sum_iW_t^i h_{\theta,t}(z_{0:t}^i),\qquad
 \mathbb E\widehat Z_t=Z_t,\qquad
 \mathbb E\widehat D_t=\nabla_\theta Z_t.
 \label{eq:disturbance-particle-derivative}
\end{equation}
\end{proposition}

\begin{proof}
At time zero, cancellation of the initial proposal gives
\[
 \mathbb E\gamma_0^N(\varphi)
 =\frac1N\sum_i\mathbb E\left[
       \frac{\rho_{0,\theta}(v_0^i)}{q_{0,\theta}(v_0^i)}
       \varphi(v_0^i)\right]
 =\int\rho_{0,\theta}\varphi\,\dd\lambda_0=\gamma_0(\varphi).
\]
For the induction step, the definitions cancel the random normalization:
\[
 \gamma_t^N(\varphi)
 =\widehat Z_{t-1}\frac1N\sum_i\omega_t^i\varphi(z_{0:t}^i).
\]
The ancestor and proposal cancellations used in
Proposition~\ref{prop:disturbance-apf-normalizer}, now with $\varphi$ in
the integrand, give the exact conditional identity
\begin{equation}
 \mathbb E[\gamma_t^N(\varphi)\mid\mathcal F_{t-1}]
 =\widehat Z_{t-1}\sum_aW_{t-1}^a
       (K_{t,\theta}\varphi)(z_{0:t-1}^a)
 =\gamma_{t-1}^N(K_{t,\theta}\varphi).
 \label{eq:disturbance-particle-induction-step}
\end{equation}
Taking expectations and applying the induction hypothesis gives
$\gamma_{t-1}(K_{t,\theta}\varphi)=\gamma_t(\varphi)$.
The argument first applies to nonnegative functions by Tonelli's theorem;
applying it to the positive and negative parts proves the signed case under
absolute integrability. Set $\varphi=1$ for the likelihood and
$\varphi=h_{\theta,t}$ componentwise for its derivative, using
Proposition~\ref{prop:disturbance-fisher} at the time prefix $t$.
\end{proof}

This proof also specifies an executable analytical recursion. A particle
need carry only its state $x_t^i$, Jacobian $\dot x_t^i$, and accumulated
score $s_t^i$, rather than an explicit array of all past disturbances.
At initialization, $s_0^i=\nabla_\theta\log\rho_{0,\theta}(v_0^i)$.
After choosing ancestor $a=a_t^i$, copy its state, Jacobian, and score;
evaluate \eqref{eq:disturbance-tangent} at the newly drawn disturbance; and,
under the differentiability assumptions in
Assumption~\ref{ass:disturbance-regularity}, set
\begin{align}
 s_t^i={}&s_{t-1}^a
 +\partial_\theta\log r_{t,\theta}(u_t^i\mid x_{t-1}^a)
 +(\dot x_{t-1}^a)^{\mathsf T}
          \nabla_x\log r_{t,\theta}(u_t^i\mid x_{t-1}^a)
 \nonumber\\
 &+\partial_\theta\log g_{t,\theta}(y_t\mid x_t^i)
 +(\dot x_t^i)^{\mathsf T}\nabla_x\log g_{t,\theta}(y_t\mid x_t^i).
 \label{eq:disturbance-analytical-score-recursion}
\end{align}
Induction on this sum gives $s_t^i=h_{\theta,t}(z_{0:t}^i)$ exactly.
Normalize the weights in \eqref{eq:disturbance-particle-measure} and report
$\widehat Z_t$, $\widehat D_t=\widehat Z_t\sum_iW_t^i s_t^i$, and
$\widehat S_t=\sum_iW_t^i s_t^i$ separately. All derivatives in this
recursion hold the realised disturbances fixed; no proposal-density
derivative enters the complete target score. Choosing
$q_{0,\theta}=\rho_{0,\theta}$, $q_{t,\theta}=r_{t,\theta}$ and
$\alpha_t=W_{t-1}$ gives the bootstrap disturbance filter, with
$\omega_0^i=1$ and $\omega_t^i=g_{t,\theta}(y_t\mid x_t^i)$.
It is the simplest comparator for guided proposals.

This forward recursion preserves the support of a degenerate transition.
Its particle histories can still coalesce under resampling, so it retains
the long-horizon variance problem of a genealogical score estimator.
Avoiding an invalid backward density resolves a support problem; it does not
prove a variance bound. The next proposition separates another issue:
unbiasedness of the unnormalized pair does not survive division in general.
Write $\widehat D_N=\widehat D_T$, $\widehat Z_N=\widehat Z_T$,
$D=\nabla_\theta Z$, and $Z=Z_T$.

\begin{proposition}[Normalization creates a separate bias question]
\label{prop:ratio-bias}
Unbiasedness of $\widehat D_N$ and $\widehat Z_N$ does not imply
$\mathbb E[\widehat D_N/\widehat Z_N]=D/Z$. In particular, if
$(\widehat Z_N,\widehat D_N)$ equals $(1/2,1)$ or $(3/2,1)$ with equal
probability, then $\mathbb E\widehat Z_N=1$, $\mathbb E\widehat D_N=1$, but
\begin{equation}
 \mathbb E\left[\frac{\widehat D_N}{\widehat Z_N}\right]
 =\frac12\left(2+\frac23\right)=\frac43\ne1
 =\frac{D}{Z}.
 \label{eq:ratio-bias-counterexample}
\end{equation}
\end{proposition}

\begin{proof}
The displayed two-point distribution has positive denominators and the stated
first moments by direct arithmetic. Its expected ratio is the average of the
two reciprocal denominators, which is $4/3$. Therefore the ratio operation
does not commute with expectation in general.
\end{proof}

The practical output of a disturbance filter should therefore retain
$\widehat Z_N$, $\widehat D_N$, and $\widehat D_N/\widehat Z_N$ as three
different quantities. A stop-gradient correction for a parameter-dependent
resampling law, such as the correction of
\citet[Section 3.1, equation (13)]{scibior2021dpf}, is relevant when automatic
differentiation is used to differentiate the random finite program. It is not
a second term to append to the Fisher estimator in
\eqref{eq:disturbance-complete-score}; that estimator already averages the
complete score under the particle importance measure. The induction proves
that no extra resampling-law term is missing from $\widehat D_t$.
Appending another term changes the estimator and requires a separate
expectation-preservation proof.

\begin{proposition}[An exact unnormalized control variate]
\label{prop:disturbance-control-variate}
In the setting of Proposition~\ref{prop:disturbance-particle-induction},
fix a parameter direction $v\in\R^p$, let $c_\theta(z)$ be a scalar,
deterministic test function with $\int\pi_\theta|c_\theta|\,\dd\lambda<\infty$,
and suppose its target integral
\begin{equation}
 C(\theta)=\int\pi_\theta(z;y)c_\theta(z)\,\dd\lambda(z)
 \label{eq:control-known-integral}
\end{equation}
is known exactly. Put $A_N=v^{\mathsf T}\widehat D_N$ and
$B_N=\gamma_T^N(c_\theta)$. For any fixed scalar $\beta$, define
\begin{equation}
 \widehat D_{\rm cv}^{\,v}=A_N-\beta(B_N-C(\theta)),\qquad
 \mathbb E\widehat D_{\rm cv}^{\,v}=v^{\mathsf T}\nabla_\theta Z.
 \label{eq:disturbance-control-variate}
\end{equation}
If $A_N,B_N$ have finite second moments and
$\operatorname{Var}(B_N)>0$, the variance-minimizing fixed coefficient is
\begin{equation}
 \beta^*=\frac{\operatorname{Cov}(A_N,B_N)}
                 {\operatorname{Var}(B_N)}.
 \label{eq:control-variate-coefficient}
\end{equation}
\end{proposition}

\begin{proof}
Proposition~\ref{prop:disturbance-particle-induction} gives
$\mathbb E A_N=v^{\mathsf T}\nabla_\theta Z$ and
$\mathbb E B_N=C(\theta)$, so
\begin{align*}
 \mathbb E[\widehat D_{\rm cv}^{\,v}]
 &=v^{\mathsf T}\nabla_\theta Z-\beta C(\theta)+\beta C(\theta)
 =v^{\mathsf T}\nabla_\theta Z.
\end{align*}
The constant does not affect variance. Its full expansion is
\begin{equation}
 \operatorname{Var}(\widehat D_{\rm cv}^{\,v})
 =\operatorname{Var}(A_N)-2\beta\operatorname{Cov}(A_N,B_N)
       +\beta^2\operatorname{Var}(B_N).
 \label{eq:control-variate-variance}
\end{equation}
Differentiating this quadratic gives
$-2\operatorname{Cov}(A_N,B_N)+2\beta\operatorname{Var}(B_N)=0$,
which proves \eqref{eq:control-variate-coefficient}; the second derivative
$2\operatorname{Var}(B_N)>0$ proves the unique minimum. At the minimum the
variance is
$\operatorname{Var}(A_N)-\operatorname{Cov}(A_N,B_N)^2/
 \operatorname{Var}(B_N)$.
\end{proof}

The proposition gives a real criterion for a control variate: its target
integral must be known, or estimated with a separately proved correction.
This is the same expectation-preserving role played by baselines in the
score-function construction of \citet[Theorem 1 and equation (4.7)]{foerster2018dice},
but the present proposition makes the required target integral explicit for
the disturbance measure. Taking $v$ as each coordinate vector gives the
components of the derivative. For one importance draw without resampling,
$A_1=R_\theta(\mathsf Z)v^{\mathsf T}h_\theta(\mathsf Z;y)$ and
$B_1=R_\theta(\mathsf Z)c_\theta(\mathsf Z)$ give the same result.
A detached KDM score has no known target integral merely because it is
empirically centered. A coefficient fitted on the same estimating particles
can destroy unbiasedness. An independent pilot fit is admissible if,
conditional on that fit, the estimating run satisfies the same identities
and the resulting products are integrable. Its estimated coefficient need
not attain the theoretical minimum.

To turn the control into a concrete construction, use a reference model whose
unnormalized disturbance density $\widetilde\pi_\theta(z;y)$ and likelihood
$\widetilde Z(\theta;y)$ are tractable. The reference must use the same
disturbance coordinates and measures. A linear-Gaussian reference can have a
singular state transition: its Gaussian observation likelihood is still
tractable when the relevant observation covariance is positive definite.
Its approximation error is corrected, rather than declared absent.

\begin{proposition}[A tractable reference supplies the centering constant]
\label{prop:disturbance-reference-control}
Suppose $\widetilde\pi_\theta$ obeys differentiation under its integral, its
likelihood $\widetilde Z$ and derivative are known exactly, and
$v^{\mathsf T}\nabla_\theta\widetilde\pi_\theta$ vanishes outside the support
of $\pi_\theta$. Define, on that support,
\begin{equation}
 c_\theta(z)=
 \frac{v^{\mathsf T}\nabla_\theta\widetilde\pi_\theta(z;y)}
      {\pi_\theta(z;y)},\qquad
 C(\theta)=v^{\mathsf T}\nabla_\theta\widetilde Z(\theta;y).
 \label{eq:disturbance-reference-control}
\end{equation}
Then this $C$ is the exact centering constant in
Proposition~\ref{prop:disturbance-control-variate}. When
$\widetilde\pi_\theta>0$, its numerator can be computed as
$\widetilde\pi_\theta v^{\mathsf T}\widetilde h_\theta$,
where $\widetilde h_\theta=\nabla_\theta\log\widetilde\pi_\theta$.
For independent draws $\mathsf Z_i\sim Q_\theta$ this gives the explicit
residual estimator
\begin{equation}
 \widehat D_{\rm ref}^{\,v}
 =\beta v^{\mathsf T}\nabla_\theta\widetilde Z+
 \frac1N\sum_{i=1}^N
 \frac{v^{\mathsf T}\nabla_\theta\pi_\theta(\mathsf Z_i;y)
       -\beta v^{\mathsf T}\nabla_\theta
                    \widetilde\pi_\theta(\mathsf Z_i;y)}
      {Q_\theta(\mathsf Z_i)}.
 \label{eq:disturbance-reference-residual}
\end{equation}
Its expectation is $v^{\mathsf T}\nabla_\theta Z$. If the reference equals
the target in a neighbourhood of $\theta$, $\beta=1$ makes this derivative
estimator deterministic and exact.
\end{proposition}

\begin{proof}
The support condition and cancellation give
\[
 \int\pi_\theta c_\theta\,\dd\lambda
 =\int v^{\mathsf T}\nabla_\theta\widetilde\pi_\theta\,\dd\lambda
 =v^{\mathsf T}\nabla_\theta\widetilde Z.
\]
The absolute-integrability requirement is precisely integrability of the
reference derivative. Substituting this control in the importance-draw
version of \eqref{eq:disturbance-control-variate} yields
\eqref{eq:disturbance-reference-residual}. Integrating its residual cancels
$Q_\theta$ and returns
$v^{\mathsf T}\nabla_\theta Z-\beta
 v^{\mathsf T}\nabla_\theta\widetilde Z$.
Adding the displayed constant proves unbiasedness. Equality of target and
reference in a neighbourhood makes their derivatives equal, so at $\beta=1$
the residual vanishes for every draw.
\end{proof}

The same control applies to the resampled filter through
$B_N=\gamma_T^N(c_\theta)$, whose unbiasedness was proved by induction.
For a factorized reference, its density ratio
$\widetilde\pi_{\theta,t}/\pi_{\theta,t}$ and complete score can be carried
along each particle's genealogy using the corresponding reference state
and tangent. Merely inserting the reference score
$\widetilde h_\theta$ omits the density ratio and is wrong relative to
\eqref{eq:disturbance-reference-control}. A poor reference can make this
ratio volatile or give an infinite second moment. The variance theorem
therefore requires its stated moment condition and a coefficient appropriate
to the complete estimator. These identities alone do not justify a
particular linearization or coefficient for a DSGE model.

Exact integration is preferable when part of the original model is
conditionally tractable. The next result states precisely what such
integration preserves and what variance comparison it permits.

\begin{proposition}[Exact conditional integration]
\label{prop:disturbance-marginalization}
Partition the disturbances as $z=(a,b)$ on a product reference measure, and
define
$\overline\pi_\theta(a)=\int\pi_\theta(a,b;y)\,\dd b$.
Assume differentiation under this conditional integral is valid and
$\overline\pi_\theta(a)>0$ on its support. Then
\begin{equation}
 \nabla_\theta\log\overline\pi_\theta(a)
 =\mathbb E_\theta[h_\theta(a,B;y)\mid a,y].
 \label{eq:disturbance-conditional-score}
\end{equation}
For a single importance draw $(A,B)$ from
$Q(a,b)=Q_A(a)Q_B(b\mid a)$ covering the target support, the scalar
unnormalized derivative $G=v^{\mathsf T}\nabla_\theta\pi_\theta(A,B;y)/
Q(A,B)$ satisfies
\begin{equation}
 \overline G(A):=\mathbb E_Q[G\mid A]
 =\frac{v^{\mathsf T}\nabla_\theta\overline\pi_\theta(A)}{Q_A(A)},
 \qquad
 \operatorname{Var}_Q(\overline G)\leq\operatorname{Var}_Q(G),
 \label{eq:disturbance-rao-blackwell}
\end{equation}
provided $G$ has finite second moment.
\end{proposition}

\begin{proof}
Differentiate the conditional integral, use
$\nabla_\theta\pi_\theta=\pi_\theta h_\theta$, and divide by
$\overline\pi_\theta$ to obtain
\eqref{eq:disturbance-conditional-score}. Conditional integration under the
proposal cancels $Q_B$:
\[
 \mathbb E_Q[G\mid A=a]
 =\frac1{Q_A(a)}
    \int v^{\mathsf T}\nabla_\theta\pi_\theta(a,b;y)\,\dd b
 =\frac{v^{\mathsf T}\nabla_\theta\overline\pi_\theta(a)}{Q_A(a)}.
\]
The law of total variance,
$\operatorname{Var}(G)=
\operatorname{Var}(\mathbb E[G\mid A])+
\mathbb E[\operatorname{Var}(G\mid A)]$, proves the inequality.
\end{proof}

For a genuinely conditionally linear-Gaussian block, a Kalman calculation
can evaluate $\overline\pi_\theta$ exactly. Replacing a nonlinear block by a
Gaussian approximation changes that integral; it belongs instead in the
corrected reference construction above. The variance inequality compares
these two coupled single-draw estimators. It does not compare complete
particle filters with different proposals and future resampling decisions.

\begin{proposition}[A moving-support Gaussian check]
\label{prop:disturbance-moving-line}
Let $U\sim\mathcal N(0,1)$, $X=(U,\theta U)^{\mathsf T}$ and
$Y=X+\varepsilon$ with independent
$\varepsilon\sim\mathcal N(0,I_2)$. Put
$d=2+\theta^2$ and $b=y_1+\theta y_2$. The observed-data score is
\begin{equation}
 \partial_\theta\log Z(\theta;y)
 =-\frac{\theta}{d}+\frac{y_2b}{d}-\frac{\theta b^2}{d^2}
 =\mathbb E[U(y_2-\theta U)\mid Y=y].
 \label{eq:disturbance-moving-line-score}
\end{equation}
\end{proposition}

\begin{proof}
The state lies on a moving line, yet its observation covariance is
$\Sigma_\theta=I_2+aa^{\mathsf T}$, $a=(1,\theta)^{\mathsf T}$.
Its determinant is $d$ and
$\Sigma_\theta^{-1}=I_2-aa^{\mathsf T}/d$. Hence
\begin{equation}
 \log Z=-\log(2\pi)-\tfrac12\log d
        -\tfrac12(y_1^2+y_2^2)+\frac{b^2}{2d}.
 \label{eq:disturbance-moving-line-likelihood}
\end{equation}
Using $d'=2\theta$ and $b'=y_2$ gives the first expression in
\eqref{eq:disturbance-moving-line-score}. Alternatively, completing the
square in
$u^2+(y_1-u)^2+(y_2-\theta u)^2$ gives
$U\mid y\sim\mathcal N(b/d,1/d)$. The complete disturbance score is
$h_\theta(u;y)=u(y_2-\theta u)$, so its posterior expectation is
$y_2b/d-\theta(1/d+b^2/d^2)$, the same expression.
\end{proof}

This example checks both the moving support and the total path derivative.
It also supplies a limiting case for the control variate: using this exact
Gaussian model as its own reference makes the $\beta=1$ derivative residual
zero. The normalized ratio with a still-random likelihood denominator need
not be exact. The likelihood can instead be evaluated analytically in this
particular example.

\begin{corollary}[What the proposed method does and does not establish]
\label{cor:disturbance-scope}
Under the assumptions above, disturbance proposals and the estimator pair
$(\widehat Z_N,\widehat D_N)$ preserve the declared degenerate model and give
an exact unnormalized likelihood/derivative identity. They do not establish a
finite-$N$ unbiased normalized score, low variance for long horizons,
global differentiability across DSGE regime switches, or an exact HMC force.
\end{corollary}

\begin{proof}
The first statement is Proposition~\ref{prop:disturbance-fisher},
Proposition~\ref{prop:disturbance-proposal}, and
Proposition~\ref{prop:disturbance-particle-induction}. The second follows from the
ratio counterexample in Proposition~\ref{prop:ratio-bias}; from the fact that
Assumption~\ref{ass:disturbance-regularity} is local and excludes unhandled
regime-boundary terms. None of these expectation identities bounds variance
as $T$ grows or equates a realised noisy score with the derivative of a
specified finite likelihood program.
\end{proof}

An exact gradient and an exact Metropolis acceptance calculation are
different requirements. For a differentiable deterministic force, momentum
and position updates in a symmetric leapfrog composition are triangular
shears with unit Jacobian determinant; the momentum flip reverses that
composition. An exact endpoint Metropolis calculation can therefore
correct such a proposal even when its force is approximate. That fact does
not make the force exact. A noisy score alone
does not specify the auxiliary-variable proposal or acceptance law needed
for a valid sampler. Those constructions and their efficiency are outside
the result proved here.


\clearpage
\bibliographystyle{plainnat}
\bibliography{ledh_younis_kdm_score}
\end{document}
